% Reciprocal proof terms and dependent constructions for the fixed2013 foundations.
\begin{annexitem}{annex:p07-hott-path-laws}{thm:p07-hott-path-laws}{Groupoid laws, cancellation and the stated pentagon}
\begin{proof}
For right unit use the J-motive $p=p\cdot\mathsf{refl}_y$ over
$(x,y,p)$; for left unit use $p=\mathsf{refl}_x\cdot p$. At reflexivity
both are witnessed by $\mathsf{refl}_{\mathsf{refl}_x}$.
The two inverse-law motives are $p^{-1}\cdot p=\mathsf{refl}_y$ and
$p\cdot p^{-1}=\mathsf{refl}_x$; double inversion has motive
$(p^{-1})^{-1}=p$. Their reflexivity values are the same specified
reflexivity of a reflexivity. For associativity take the motive over
all three composable paths and their endpoints. Successive J on $p,q,r$
reduces both sides to $\mathsf{refl}_x$, so defines the associator.
No induction with two fixed endpoints was used.

Left composition by $p:x=y$ has inverse left composition by $p^{-1}$:
the two composites reduce by the associator, inverse law and unit law
to the original argument. Right composition is treated with the same
three laws in reverse-side order. These explicit inverse homotopies
justify cancellation. The five parenthesizations of four composable
paths have the canonical J-associators between them. The pentagon's
two composites are paths with the same source and target; whisker its
associators using ap of the relevant concatenation functions. J on all
four original paths, with the endpoints varying, reduces every associator
and whiskered associator to reflexivity. Both pentagon routes reduce
to that reflexivity by the stated double-J computation. Reflexivity
therefore defines their comparison, a path between the two higher paths.
This proves that particular pentagon, not an unstated entire operad.
\end{proof}
\end{annexitem}

\begin{annexitem}{annex:p07-hott-sigma-paths}{thm:p07-hott-sigma-paths}{Pair encoding, decoding and transport}
\begin{proof}
Encode a total-space path by J, taking reflexivity at $w$ to
$(\mathsf{refl}_{\pi_1w},\mathsf{refl}_{\pi_2w})$.
For decoding, eliminate both endpoint pairs and the input pair $(p,q)$.
Based J on $p:x=y$ reduces $y$ to $x$ and transport to its input;
J on $q:u=v$ then reduces $v$ to $u$. Return reflexivity of $(x,u)$.
For decode after encode, J on the original total-space path followed
by pair elimination makes the composite reflexivity of the same pair.
For encode after decode, eliminate the pairs, then J on $p$ and $q$;
the composite returns the same pair of reflexivities. These two J-built
comparisons are the quasi-inverse homotopies; adjointification turns
them into the stated equivalence. This also proves pair uniqueness
by encoding the pair of reflexivities, and gives the specified lift
by decoding $(p,\mathsf{refl}_{\mathsf{transport}(p,u)})$.

For product transport use J on the base path, then pair elimination
on the input: both sides compute to that pair. For dependent-sum
transport use the same order. At reflexivity the first transport is
$u$, its lift is reflexivity of $(x,u)$, and the second transport is
$z$; hence both sides are $(u,z)$. Applying J yields the asserted
path in the target total space. A path between pairs sharing a first
coordinate still encodes an arbitrary loop at that coordinate; the
proof never replaces it by reflexivity without justification.
\end{proof}
\end{annexitem}

\begin{annexitem}{annex:p07-hott-function-paths}{thm:p07-hott-function-paths}{The explicitly assumed inverse comparisons}
\begin{proof}
Take the inverse function and its two homotopies from the assumed
equivalence of happly. The counit is a path between pointwise homotopies;
evaluating it at $x$ yields the stated beta path. The unit is the
path eta comparison directly. J on function paths shows that happly
sends an inverse or a composite to pointwise inverses or composites,
since all base paths become reflexivity. Its inverse comparisons
then identify the original function paths with their funext images.

For arrow transport apply J to $p$. Its right-hand expression reduces
to $\lambda v.f(v)$ and its left to $f$; source Pi-eta supplies equality.
For dependent functions the precise right-hand value at $v:P(y)$ is
\[
 \mathsf{transport}^{Q}
 (\mathsf{lift}(v,p^{-1})^{-1},
                       f(\mathsf{transport}^{P}(p^{-1},v))).
\]
The lift runs from $(y,v)$ to $(x,\mathsf{transport}^{P}(p^{-1},v))$,
so its inverse transports the result forward in the correct family.
J on $p$ reduces this pointwise formula to $f(v)=f(v)$; funext makes
it a path of dependent functions where required. For the dependent-path
equivalence, J on $p$ reduces the displayed source and target to
$f=g$ and $\Pi_u(f(u)=g(u))$. The specified happly equivalence and
its inverse give both directions and both round trips. Their
propositional beta comparison gives the pointwise computation; this
also treats the nondependent family specialization.
\end{proof}
\end{annexitem}

\begin{annexitem}{annex:p07-hott-contractibility-laws}{thm:p07-hott-contractibility-laws}{All stated contractions and removals}
\begin{proof}
For $c:\Pi_x(a=x)$, apd at $p:x=y$ and the equality-family transport
formula give $c_x\cdot p=c_y$. Cancelling yields
$p=c_x^{-1}\cdot c_y$. That composite is a center of the path type;
inverting this comparison contracts every path to it. In particular
$c_a=\mathsf{refl}_a$, since every loop at $a$ is this composite.
A contraction makes $A$ inhabited and supplies $x=y$ by this canonical
path. Conversely an inhabitant $a$ and a proposition witness supply
the contraction $x\mapsto(a=x)$. The maps $A\to\mathbf1$ and
$\mathbf1\to A$, the latter constant at $a$, have inverse homotopies
from the contraction and unit induction. In the other direction an
equivalence with unit transports its contraction along the inverse
comparisons, or along the following retract construction.

For a retract $r:A\to B$ with section $s$ and $\epsilon:r s\sim\mathsf{id}$,
use center $r(a)$ and contraction
$\mathsf{ap}_r(c_{s(b)})\cdot\epsilon_b:r(a)=b$.
The based-path total space has center $(a,\mathsf{refl}_a)$:
J over $(x,p:a=x)$ gives its path to $(x,p)$, computing to reflexivity
at the center. If each $P(x)$ has center $d_x$, the first projection
$\Sigma_xP(x)\to A$ has inverse $x\mapsto(x,d_x)$; one composite
is judgmentally the identity, and the other has its pair-path comparison
from each fiber contraction over reflexivity.

If $A$ has contraction $c$, the map to $P(a)$ sends $(x,u)$ to
$\mathsf{transport}^{P}(c_x^{-1},u)$; its inverse sends $v$ to $(a,v)$.
The first round trip uses $c_a=\mathsf{refl}_a$. The other uses the
pair path with base $c_x$ and inverse-transport comparison
$\mathsf{transport}(c_x,\mathsf{transport}(c_x^{-1},u))=u$.
For a contractible base and contractible fibers, center $(a,d_a)$
contracts to $(x,u)$ by the base contraction and the target-fiber path
from its transported second component to $u$.

Under funext, pointwise centers form a dependent function, and the
pointwise contractions become a contraction of that function space.
For two contractions $(a,c),(a',c')$ of one type, use $c_{a'}:a=a'$.
After transport both second components have value in the same path
type $a'=x$. That path type is contractible by the first paragraph,
so their values are equal. Funext gives equality of the second
components and the pair-path equivalence gives equality of the
contractions. Thus isContr is a mere proposition. All constructions
produce the specified data; no mere-existence choice has been used.
\end{proof}
\end{annexitem}

\begin{annexitem}{annex:p07-hott-prop-set}{thm:p07-hott-prop-set}{Propositions and all stated closure countercases}
\begin{proof}
Given a proposition witness and arbitrary $x:A$, it makes $A$
contractible with center $x$. Hence every $x=y$ is contractible,
and any two paths in it are equal. This proves that a proposition
is a set, including the empty case since $x$ was supplied as an
argument, not chosen globally. If $P,Q$ are mere propositions and
maps run in both directions, both composite maps are pointwise
equal to the identity by their proposition witnesses; adjointification
therefore supplies the equivalence.

Under funext, two witnesses of isProp agree at each pair $x,y$,
because $A$ is then a set. Repeated funext identifies the witnesses.
For isSet, each identity type is a proposition and hence a set;
its parallel higher paths are equal, so the witness functions again
agree pointwise and repeated funext identifies them. A dependent
product of propositions is a proposition by applying funext to
the supplied pointwise equality of any two sections.

For a proposition-valued $P$, the pair-path encoding has fiber
$\mathsf{transport}(p,u)=v$ over each base path $p$. That fiber is
contractible because $P(y)$ is a proposition. Removing these
contractible fibers proves that first projection induces an
equivalence on paths. Product propositionhood follows by pairing
the two coordinate paths. For the countercase, in $\mathbf2=\mathbf1+\mathbf1$
the family taking the left constructor to $\mathbf1$ and the right
to $\mathbf0$ turns an alleged path between constructors into an
inhabitant of $\mathbf0$ by transport. The constructors are unequal,
so Bool is not a proposition. If $\Sigma_{x:\mathbf2}\mathbf1$ were
a proposition, comparing $(0,\star)$ and $(1,\star)$ and applying
first projection would give that forbidden Boolean path. Thus
proposition-valued fibers alone do not make a total space a proposition.
\end{proof}
\end{annexitem}

\begin{annexitem}{annex:p07-hott-transport-laws}{thm:p07-hott-transport-laws}{All actions and transported identity formulas}
\begin{proof}
J with reflexivity value $u$ defines transport; with reflexivity value
$\mathsf{refl}_{f(x)}$ it defines ap; the motive
$\mathsf{transport}^{P}(p,s(x))=s(y)$ defines apd with reflexivity value
$\mathsf{refl}_{s(x)}$. Their computations are thus the stated ones.
For preservation of concatenation apply J to both composable paths;
for inverse preservation apply J to the one path. For action of a
composite, and for action of the identity function, apply J to the path.
Each resulting base is equality of two identical reflexivity paths.

Transport composition uses J on $p$ and $q$, reducing both sides to $u$.
For a pulled-back family, J on $p$ reduces both transports to $u$.
The fiber-map comparison reduces after J to $k_x(u)=k_x(u)$.
Transport along $p^{-1}$ is inverse transport because composition,
inverse paths and reflexivity computation reduce each composite to its
input. For a constant family J gives the comparison
$c(p,u):\mathsf{transport}(p,u)=u$. J on $p$ proves
$\mathsf{apd}_f(p)=c(p,f(x))\cdot\mathsf{ap}_f(p)$, since all three
paths become reflexivity in that case. Composing with $c$ or $c^{-1}$
gives inverse maps between the ordinary and dependent path descriptions.

For the three nondependent equality families, use J on $p$. The
remaining equations reduce respectively to $r=r\cdot\mathsf{refl}$,
$r=\mathsf{refl}\cdot r$, and
$r=\mathsf{refl}\cdot r\cdot\mathsf{refl}$ after ap of an identity
or constant function is evaluated. The unit laws provide these paths.
For dependent $f,g$, the middle map acts on $r$ between the transported
endpoints. At reflexivity its action is ap of the identity on that
fiber, while both apd paths are reflexivities. The same ap-identity
and unit paths prove the base. This gives the corrected typed dependent
formula without equating elements from different fibers.
\end{proof}
\end{annexitem}

\begin{annexitem}{annex:p07-hott-homotopy-natural}{thm:p07-hott-homotopy-natural}{The naturality square and cancellation}
\begin{proof}
Apply J to $p$ with the displayed naturality square as its motive.
At reflexivity the desired comparison is
$H(x)\cdot\mathsf{refl}_{g(x)}=\mathsf{refl}_{f(x)}\cdot H(x)$;
the two unit paths identify both sides with $H(x)$.
For $H:k\sim\mathsf{id}$, apply naturality at the path
$H(x):k(x)=x$. It gives
$H(k(x))\cdot H(x)=\mathsf{ap}_k(H(x))\cdot H(x)$,
using ap of the identity. Right cancellation proved above yields the
self-composite comparison. Reflexivity, inverse and concatenation of
homotopies are the pointwise constructions, so no equality of functions
or function-extensionality premise was silently used.
\end{proof}
\end{annexitem}

\begin{annexitem}{annex:p07-hott-eckmann-hilton}{thm:p07-hott-eckmann-hilton}{Typed horizontal operations on second paths}
\begin{proof}
For $\alpha:p=q$ in $a=b$ and $\beta:r=s$ in $b=c$, define
\[
 \alpha\star\beta=
 \mathsf{ap}_{t\mapsto t\cdot r}(\alpha)\cdot
 \mathsf{ap}_{t\mapsto q\cdot t}(\beta).
\]
The alternative first changes $r$ to $s$, then $p$ to $q$:
\[
 \alpha\star'\beta=
 \mathsf{ap}_{t\mapsto p\cdot t}(\beta)\cdot
 \mathsf{ap}_{t\mapsto t\cdot s}(\alpha).
\]
Both paths run from $p\cdot r$ to $q\cdot s$. With all these endpoints
free in the motive, J on $\alpha$ and $\beta$ reduces both composites
to reflexivity, giving $\star=\star'$ pointwise.

Now specialize $p,q,r,s$ to $\mathsf{refl}_a$. Right concatenation
with this unit acts as the identity on second loops: the canonical
right-unit homotopy on the path type has reflexivity as its value at
$\mathsf{refl}_a$, so its naturality square at a second loop, followed
by the two unit cancellations, identifies this action with that loop.
The canonical left-unit homotopy gives the same assertion for left
concatenation. Consequently $\alpha\star\beta=\alpha\cdot\beta$
and $\alpha\star'\beta=\beta\cdot\alpha$. The previously constructed
comparison yields the desired commutativity. This specialization uses
typed unit coherences; it does not perform illegal J on arbitrary loops
while keeping both endpoints fixed.
\end{proof}
\end{annexitem}

\begin{annexitem}{annex:p07-hott-adjointification}{thm:p07-hott-adjointification}{The actual adjusted counit and its triangle}
\begin{proof}
The revised composite has endpoints
$f g b\to f g f g b\to f g b\to b$, so is typed.
Put $t=\mathsf{ap}_f(\eta_a)$. The self-composite law for
$\eta:g f\sim\mathsf{id}$ gives
$\eta_{g f a}=\mathsf{ap}_g(t)$.
Naturality of $\epsilon:f g\sim\mathsf{id}$ at
$t:f g f a=f a$ gives
\[
 \epsilon_{f g f a}\cdot t
 =\mathsf{ap}_{f g}(t)\cdot\epsilon_{f a}.
\]
Functoriality and this square identify
$\epsilon'_{f a}$ with
$\epsilon_{f g f a}^{-1}\cdot\epsilon_{f g f a}\cdot t=t$.
Invert that comparison to obtain the required
$\tau_a:t=\epsilon'_{f a}$. Thus a genuine triangle is constructed,
not stipulated for the unadjusted quasi-inverse. Forgetting a supplied
triangle gives quasi-inverse data immediately.

For a left inverse $(l,\eta)$ and right inverse $(r,\epsilon)$,
construct $\gamma_y:r(y)=l(y)$ as
$\eta_{r(y)}^{-1}\cdot\mathsf{ap}_l(\epsilon_y)$.
Then $(r,\lambda x.\gamma_{f(x)}\cdot\eta_x,\epsilon)$ is a
quasi-inverse of $f$. Conversely a quasi-inverse supplies both halves
of bi-invertibility. These maps join the data explicitly; inverse
functions were never assumed judgmentally identical.
\end{proof}
\end{annexitem}

\begin{annexitem}{annex:p07-hott-fiber-equivalence}{thm:p07-hott-fiber-equivalence}{Typed fiber contractions in both directions}
\begin{proof}
The Sigma-path equivalence first gives
$\Sigma_\gamma(\mathsf{transport}^{x\mapsto(f(x)=b)}(\gamma,p)=p')$.
Equality-family transport is $\mathsf{ap}_f(\gamma)^{-1}\cdot p$;
composing with $\mathsf{ap}_f(\gamma)$ converts its comparison to
$\mathsf{ap}_f(\gamma)\cdot p'=p$, with inverse conversion by the
inverse path. The groupoid laws prove both conversions are inverse.

For half-adjoint data use center $(g b,\epsilon_b)$ and
$\gamma=\mathsf{ap}_g(p)^{-1}\cdot\eta_a:g b=a$.
The required second component is supplied by
\[
 \mathsf{ap}_f(\gamma)\cdot p
 =\mathsf{ap}_{f g}(p)^{-1}\cdot\epsilon_{f a}\cdot p
 =\mathsf{ap}_{f g}(p)^{-1}\cdot\mathsf{ap}_{f g}(p)\cdot\epsilon_b
 =\epsilon_b.
\]
The first equality uses the triangle, the second naturality of
$\epsilon$, and the third inverse/unit comparisons. Decoding the
resulting fiber pair gives its contraction.

Conversely write the specified center of each fiber as $(g b,\epsilon_b)$.
At $b=f(a)$ apply its contraction to $(a,\mathsf{refl}_{f a})$.
Encode that fiber path. Its first component is $\eta_a:g f a=a$;
its second is $\mathsf{ap}_f(\eta_a)\cdot\mathsf{refl}_{f a}=\epsilon_{f a}$.
The right-unit comparison supplies $\tau_a$. These are actual
dependent functions of $a,b$, extracted from the given contractions.
No extra coherence or uniqueness axiom was inserted. A quasi-inverse
therefore yields fiber contractions by adjointifying it first.
\end{proof}
\end{annexitem}

\begin{annexitem}{annex:p07-hott-equivalence-bridges}{thm:p07-hott-equivalence-bridges}{Propositionhood without erasing inverse data}
\begin{proof}
Assume funext. For a map with inverse data $(g,\eta,\epsilon)$,
postcomposition by $f$ on $B\to A$ has inverse postcomposition by
$g$. The two composites are pointwise identified with the identity
using $\epsilon$ or $\eta$ at the value of the argument function;
funext turns these into the required function paths. Precomposition
by $f$ from $B\to A$ to $A\to A$ has inverse precomposition by $g$;
its comparisons use ap of the argument function on $\eta$ and
$\epsilon$, followed by funext. Adjointification and the fiber
criterion make the fibers of both composition maps contractible.

The fiber of postcomposition over $\mathsf{id}_B$ is equivalent to
$\mathsf{rinv}(f)$, and that of precomposition over $\mathsf{id}_A$
to $\mathsf{linv}(f)$: convert the second component using the specified
happly/funext equivalence, leaving the inverse function fixed.
Pair elimination and the two funext inverse comparisons prove both
round trips of these total-space conversions. Hence both inverse
types are contractible, not merely inhabited.

Given half-adjoint data once, fix any right inverse $(h,\delta)$.
The remaining type of $(\eta,\tau)$ is equivalent to
\[
 \Pi_{a:A}\big((h f a,\delta_{f a})=(a,\mathsf{refl}_{f a})
                                  \text{ in }\mathsf{fib}_f(f a)\big).
\]
Indeed evaluate $(\eta,\tau)$ at each $a$ and decode by the fiber
path equivalence; conversely encode each path and assemble its two
components. The pointwise round trips follow from that equivalence,
and funext and Sigma-path encoding give the total round trips.
Each displayed path type is contractible because its fiber is
contractible. Their dependent product is contractible. Reassociating
half-adjoint data as a right inverse followed by these remaining
coherent data therefore gives a contractible base with contractible
fibers. Thus ishae is contractible whenever inhabited, and is a
mere proposition for arbitrary $f$.

If biinv is inhabited, its explicit quasi-inverse gives the same
contraction of left and right inverse types; their product is
contractible. Thus biinv too is a mere proposition. The map-level
isContr is a dependent product of proposition-valued type-level
isContr predicates, so is a mere proposition. The mutual maps already
constructed between ishae and fiber contraction, and between biinv
and ishae through quasi-inverses, have inverse homotopies because
their codomains are mere propositions. Adjointification gives the
pairwise type equivalences. No propositionhood of raw qinv was used.

Identity is its own quasi-inverse. Inverting an equivalence exchanges
its inverse function with the original and the two inverse homotopies.
For $f:A\to B$ and $h:B\to C$ with inverses $g,k$, use $g k$ as
inverse of $h f$: one comparison follows from ap of $g$ on the
$k h$ comparison and then $g f\sim\mathsf{id}$; the other uses ap
of $h$ on $f g\sim\mathsf{id}$ and then $h k\sim\mathsf{id}$.
Their pointwise composites are typed homotopies. Adjointification
supplies the corresponding equivalence evidence.
\end{proof}
\end{annexitem}

\begin{annexitem}{annex:p07-hott-path-equivalence}{thm:p07-hott-path-equivalence}{Both path-map round trips with the triangle}
\begin{proof}
For $p:a=a'$, apply naturality of $\eta:g f\sim\mathsf{id}$:
$\eta_a\cdot p=\mathsf{ap}_{g f}(p)\cdot\eta_{a'}$.
Left cancellation identifies
$\eta_a^{-1}\cdot\mathsf{ap}_g(\mathsf{ap}_f(p))\cdot\eta_{a'}$
with $p$. For the opposite round trip, apply ap-functoriality to
the proposed inverse on $q:f a=f a'$. The triangle changes the result to
$\epsilon_{f a}^{-1}\cdot\mathsf{ap}_{f g}(q)\cdot\epsilon_{f a'}$.
Naturality of $\epsilon$ identifies this with $q$. These are the
two quasi-inverse homotopies, which adjointification upgrades to
equivalence data. The corrected triangle, not arbitrary raw inverse
homotopies, was used in the second calculation.
\end{proof}
\end{annexitem}

\begin{annexitem}{annex:p07-hott-level-closure}{thm:p07-hott-level-closure}{Induction on every claimed level}
\begin{proof}
For cumulativity at $n=-2$, the canonical-path contraction proves
every identity type contractible. For the successor step apply
the induction hypothesis to each identity type.
For retract closure the contractible base is already proved.
Suppose $r:X\to Y$ has section $s$ and $\epsilon:r s\sim\mathsf{id}$.
At successor level, the map ap of $s$ embeds the path type $y=y'$
as a retract of $s(y)=s(y')$, with retraction
$q\mapsto\epsilon_y^{-1}\cdot\mathsf{ap}_r(q)\cdot\epsilon_{y'}$.
The naturality square for $\epsilon$ proves its composite with
$\mathsf{ap}_s$ is the identity. The lower-level induction hypothesis
therefore applies to these path types. An equivalence supplies such
a retract in both directions, so level preservation follows.

For Sigma closure, the base $-2$ uses the explicit contraction of a
contractible total space. At the next level, its path type is
equivalent to $\Sigma_{p:a=a'}(\mathsf{transport}(p,u)=u')$.
The base and every fiber in this expression have the lower level
by the hypotheses; the Sigma induction hypothesis and equivalence
preservation give that level for the path type, hence the claimed
successor level for the total space.

For Pi closure, the contractible base uses pointwise centers and
funext. At successor level, funext identifies a function path type
with $\Pi_a(f(a)=g(a))$. Every fiber has the lower level, so the
Pi induction hypothesis and equivalence preservation give the result.
Finally isContr is a mere proposition at the base; at each successor,
the level predicate is a dependent product of lower-level predicates.
Those are mere propositions by induction, and their product is one
by funext. This proves every stated level and its predicate standing.
\end{proof}
\end{annexitem}

\begin{annexitem}{annex:p07-hott-univalence-comparison}{thm:p07-hott-univalence-comparison}{The particular inverse and the Boolean loop}
\begin{proof}
The counit $\mathsf{idtoeqv}(\mathsf{ua}(e))=e$, projected to functions
and evaluated, gives transport computation. J computes idtoeqv at
reflexivity to the identity function; propositionhood of equivalence
evidence identifies it with the identity equivalence. The unit then
gives $\mathsf{ua}(\mathsf{id})=\mathsf{refl}$.
Transport composition identifies the underlying function of idtoeqv
at $p\cdot q$ with the composite in traversal order. Funext and
propositionhood of the evidence give equality of the equivalences;
applying ua and both inverse comparisons proves the composition formula.
Transport along $p^{-1}$ has both inverse homotopies for transport
along $p$. Two inverse functions are homotopic by the explicit
left/right inverse joining construction; funext and propositionhood
of their evidence identify the equivalences. Applying ua gives the
inverse formula. Thus none of these laws is judgmental ua computation.
Boolean elimination defines interchange $e(0)=1,e(1)=0$ and in both
cases proves $e(e(x))=x$. Adjointification gives its equivalence data.
If $\mathsf{ua}(e)=\mathsf{refl}$, applying idtoeqv and evaluating
at $0$ would give $1=0$, contradicting Boolean disjointness.
Hence the universe is not a set.
\end{proof}
\end{annexitem}

\begin{annexitem}{annex:p07-hott-truncation-boundary}{thm:p07-hott-truncation-boundary}{Allowed elimination and uniform-selection obstruction}
\begin{proof}
The truncation recursor on $\mathsf{id}_P$ gives $\|P\|\to P$.
The constructor gives the other map; both inverse homotopies follow
from propositionhood of their types. Adjointification yields the
equivalence. Applying these recursors pointwise gives unique choice.
Suppose $\sigma:\Pi_A(\|A\|\to A)$ were supplied. Put
$p=\mathsf{ua}(e)$ for Boolean interchange and $u=|0|$.
Dependent action of $\sigma$ along $p$, the arrow-transport formula
and evaluation give
\[
 \mathsf{transport}^{A\mapsto A}
 (p,\sigma_{\mathbf2}(\mathsf{transport}^{A\mapsto\|A\|}(p^{-1},u)))
 =\sigma_{\mathbf2}(u).
\]
Its transported input equals $u$ because $\|\mathbf2\|$ is a
proposition. Univalence computation thus gives
$e(\sigma_{\mathbf2}(u))=\sigma_{\mathbf2}(u)$, impossible in either
Boolean constructor case. For uniform double-negation elimination
repeat this argument with the proposition-valued input family
$\neg\neg A$ and $u(w)=w(0):\neg\neg\mathbf2$. It is a proposition
by funext and propositionhood of the empty type, so the transported
input is again $u$ and the same contradiction follows.
An inhabitant of $\Pi_A(A+\neg A)$ would give that elimination:
on $v:\neg\neg A$ its left branch supplies $a$, while its right
branch supplies $w:\neg A$ and eliminates $v(w):\mathbf0$.
This concerns uniform information-bearing operations on all types,
not excluded middle restricted to mere propositions.
\end{proof}
\end{annexitem}

\begin{annexitem}{annex:p07-hott-quasi-inverse-countercase}{thm:p07-hott-quasi-inverse-countercase}{A global loop family, not an isolated loop}
\begin{proof}
Funext identifies qinv of identity with
$\Sigma_g((g=\mathsf{id})\times(g=\mathsf{id}))$.
Reassociate over the contractible based-path space $\Sigma_g(g=\mathsf{id})$;
removing it leaves $\mathsf{id}=\mathsf{id}$, which funext identifies
with $\Pi_a(a=a)$. Pairing/unpairing and Sigma elimination give both
reassociation round trips. This first equivalence needs no univalence.

Assume now univalence and truncation, and put
$x_0=(\mathbf2,|\mathsf{refl}|):X$. Subtype paths followed by
univalence identify paths in $X$ with equivalences of its carriers.
Bool is a set: fix a constructor and take code equal to unit there
and empty at the other. Encode by transporting the unit point;
decode returns reflexivity in the same-constructor branch and uses
empty elimination in the other. J gives decode-after-encode, and
Boolean/unit elimination gives encode-after-decode. Thus each
Boolean path type is unit or empty, hence a proposition.
An equivalence of Bool is identity or interchange: its value at
$0$ chooses the case, and its value at $1$ must differ, since ap of
a quasi-inverse would otherwise identify the two constructors.
Boolean elimination determines both values; funext and propositionhood
of evidence identify the corresponding equivalence. Evaluation at
$0$ separates them. The equivalence type is also a set by Pi/Sigma
closure. Every carrier in $X$ is a set by eliminating its merely
given Boolean equality into isSet, a mere proposition. Hence every
path type of $X$ is a set, so $X$ is a 1-type.

Let $q:x_0=x_0$ correspond to interchange. It is not reflexivity;
every other loop is identity or $q$, so $q$ commutes with every loop,
by the unit laws or reflexivity of $q\cdot q$. Each $x:X$ has
$\|x_0=x\|$, by its defining truncated universe path and the
subtype-path construction. Set
\[
 B(x)=\Sigma_{r:x=x}\Pi_{s:x_0=x}(r=s^{-1}\cdot q\cdot s).
\]
To show $B(x)$ a proposition, eliminate $\|x_0=x\|$ into the
mere proposition isProp$(B(x))$ and take $p:x_0=x$. Two such pairs
have equal first components by evaluating their second components
at $p$; their transported second components agree by funext since
they are paths in the set $x=x$. Sigma encoding identifies the pairs.
Now eliminate $\|x_0=x\|$ into $B(x)$ itself. On a chosen $p$
take $r=p^{-1}\cdot q\cdot p$. For any $s$, commutation of $q$
with $p\cdot s^{-1}$ and inverse cancellation give the required
$r=s^{-1}\cdot q\cdot s$. First projection yields a global
$h:\Pi_x(x=x)$; its second component at $x_0$ evaluated at
reflexivity gives $h(x_0)=q$. Equality with the reflexivity section
would evaluate to $q=\mathsf{refl}$, a contradiction. Thus the
section type has two unequal points and is not a proposition;
the first equivalence transfers this to qinv$(\mathsf{id}_X)$.
$X$ lies in a larger universe than its quantified carriers.
\end{proof}
\end{annexitem}
